\scanpage{152}
\unsure{,|1} tel qu'on ait
\[
 \boxed{\left\{\begin{aligned}
 f(x_1+\alpha_1,x_2)&=f(x_1,x_2)+f^0(\alpha_1,\sigma_2x_2-x_2)\\
 f(x_1,x_2+\alpha_2)&=f(x_1,x_2)+f^0(\sigma_1x_1-x_1,\alpha_2)
 \end{aligned}\right.}\tag*{$(*)$}
\]
$\forall\,\alpha_1\in A_1^0,\ \alpha_2\in A_2^0$. On en conclut
\[
 \boxed{\begin{aligned}
 f(x_1+\alpha_1,x_2+\alpha_2)&=f(x_1,x_2)+f^0(\alpha_1,\sigma_2x_2-x_2)+f^0(\sigma_1x_1-x_1,\alpha_1)\\
 &\hspace{8em}-2f^0.(\alpha_1,\alpha_2)
 \end{aligned}}
\]
Si on choisit $x_1,x_2$ comme origines, pour identifier $A_1$ à $A_1^0$ (via les $\alpha_1$) et $A_2$ à $A_{\unsure{3|2}}^0$ (via les $\alpha_2$), posant
\[
 \begin{aligned}
 \sigma_1x_1-x_1&=t_1-2x_1=\tau_1\in A_1^0\\
 \sigma_2x_2-x_2&=t_2-2x_2=\tau_2\in A_2^0
 \end{aligned}
\]
la formule devient
\[
 \underset{\substack{\Vert\,\mathrm{dfn}\\ f(x_1+\alpha_1,x_2+\alpha_2)}}{g(\alpha_1,\alpha_2)}
 =\deleted{\unsure{f(x}}\underset{\substack{\Vert\\ f(x_1,x_2)}}{c}+f^0(\alpha_1,\tau_2)+f^0(\tau_1,\alpha_2)-2f^0(\alpha_1,\alpha_2)
\]
% slanted
\margin{Si $2.1_k=0$, alors $\tau_1=t_1$, $\tau_2=t_2$ sont indép. du choix de $x_1,x_2$, la formule se réduit à $g(\alpha_1,\alpha_2)=c+f^0(\alpha_1,\tau_2)+f^0(\tau_1,\alpha_2)$.}

\deleted{\unsure{Inversement},} Cette formule définit \unsure{entièremt|entièrement} la fonction $g$ \inserted{donc $f$} quand on connaît $c$ ($=f(x_1,x_2)=g(0,0)$), \deleted{\unsure{et}} $\tau_1\in A_1^0,\tau_2\in A_2^0$ (qui définissent les involutions sur $A_1\simeq A_1^0$ \inserted{via $x_1$} et $A_2\simeq A_2^0$ \inserted{via $x_2$} par $\begin{aligned}\alpha_1&\longmapsto\tau_1-\alpha_1\\\alpha_2&\longmapsto\tau_2-\alpha_2\end{aligned}$) et l'application bilin. $f^0$. Quand celle-ci est donnée, on vérifie que la formule précédente définit bien une application $(\sigma_1,\sigma_2)$-bil. Mais si on ne suppose pas $2$ inv., il n'y a aucune
% slanted
\margin{Si $2$ inv., \deleted{\unsure{on}} $\tau_1=0,\tau_2=0$ (en choisissant $x_1,x_2$ fixes par \unsure{$\tau_1,\tau_2$|$\sigma_1,\sigma_2$}), \inserted{via $A_i^0$}, donc $g(\alpha_1,\alpha_2)=-2f^0(\alpha_1,\alpha_2)$ la condition est équiv. à la bilinéarité, sauf sur la const.}
% archive sheet number: 151

